\documentclass{article}
\usepackage[a4paper, total={6in, 8in}]{geometry}
\usepackage{graphicx} % Required for inserting images
\usepackage{amsmath,amssymb,xcolor,bm,amsthm,mathtools,comment, enumerate}
\newcommand{\eps}{\varepsilon}
\usepackage{tikz}
\newtheorem{theorem}{Theorem}[section]
\newtheorem{proposition}[theorem]{Proposition}
\newtheorem{remark}[theorem]{Remark}
\newtheorem{definition}[theorem]{Definition}
\newtheorem{lemma}[theorem]{Lemma}
\newtheorem{corollary}[theorem]{Corollary}
\newtheorem{example}[theorem]{Example}
\newtheorem{assumption}[theorem]{Assumption}
\newtheorem{conjecture}[theorem]{Conjecture}
\newtheorem{question}{Question}

\usepackage{cleveref}
\crefname{originaltheorem}{original theorem}{original theorems}
\Crefname{originaltheorem}{Original theorem}{Original theorems}
\makeatletter
\newlabel{t:perfect}{{1.2}{source}{}{sourcepointer-1}{}}
\newlabel{t:perfect@cref}{{[originaltheorem][0][]1.2}{[1][0][]source}{}{}{}}
\makeatother
\title{A perfect even tuple dominating the halved coordinates of a bounded even tuple}
\author{Joshua Erde}
\date{arxiv:2404.03950v3; CC BY 4.0}
\begin{document}
\maketitle
\noindent\textit{Editorial note.} The counted unit is the article's own Lemma l:first at source lines 154-156 together with the proof environment that begins on the next line (157-170); the pairing is the ordinary adjacent one and the proof ends before the author's proof of Theorem t:perfect opens at source line 172.\par\medskip
\noindent\textit{Editorial note.} The retained context at source lines 30-36 and 51 is the article's own statement of the Balister-Gyori-Schelp necessary condition and conjecture and the article's own definitions of profile, admissible, even and perfect tuples.\par\medskip
\noindent\textit{Editorial note.} The retained context at source line 64 is the article's own definition of the precedence order and of the concatenation of tuples; it is retained because the counted lemma concludes y <= y' and because the down-closedness of admissible tuples stated there is what makes the dominating tuple useful.\par\medskip
\noindent\textit{Editorial note.} The retained context at source line 152 is the article's own definition of the quantity o(x); it is retained because the hypothesis x\_\{n+1\} >= 2 o(x) of the counted lemma is stated in terms of it.\par\medskip
\noindent\textit{Editorial note.} The native preamble is retained from source lines 2-18, that is the article class, geometry, AMS, bm, mathtools, comment, enumerate and tikz packages and the author's theorem declarations; the title and author lines at source lines 19-20 are not part of the preamble fragment and the document body is not retained.\par\medskip
\noindent\textit{Editorial note.} No retained line contains a citation, so the material carries no bibliography; the retained lines refer only to the label c:BGS, which the retained conjecture at source lines 34-36 defines, and to the label t:perfect, which is left undefined by the retained fragments and is therefore declared as a non-load-bearing navigation pointer rather than as a premise.\par\medskip
\noindent\textit{Editorial note.} The article's Theorem t:perfect, Lemma l:admissible and discussion theorem c:admissible are not counted here and are not used as premises.\par\medskip
\noindent\textit{Allowed standard prerequisites (editorial).}\begin{itemize}
\item Profiles of matchings in the hypercube Q\textasciicircum{}n. For a matching M in the cube Q\textasciicircum{}n the profile is the vector x in N\textasciicircum{}n whose i-th coordinate counts the edges of M whose endpoints differ in the i-th coordinate; a tuple is admissible when it is the profile of some matching, even when every coordinate is even, and perfect when the sum of the coordinates equals 2\textasciicircum{}\{n-1\}; the article's own definitions are retained at source line 51.. Fixes the meaning of admissible, even and perfect for the counted lemma, whose conclusion asserts the existence of a perfect even tuple in N\textasciicircum{}n.
\item Precedence order on tuples and down-closedness of admissible tuples. For x and z in N\textasciicircum{}n the article writes z precedes x, in symbols z <= x, when z\_i <= x\_i for every i; the set of admissible tuples is down-closed in this partial order, and (x,z) denotes the concatenation of x in N\textasciicircum{}n and z in N\textasciicircum{}m as a vector in N\textasciicircum{}\{n+m\}; retained at source line 64.. The counted lemma concludes y <= y' for a suitable perfect even y' in N\textasciicircum{}n, which is exactly this order; the down-closedness is why a perfect even tuple dominating y transfers admissibility to y.
\item The parity obstruction of Balister, Gyori and Schelp. If the cube F\_2\textasciicircum{}n is partitioned into pairs \{a\_i,b\_i\} with differences d\_i, then the sum of all d\_i vanishes, because each vertex occurs exactly once on each side of the sum; the article records this identity at source lines 30-33 and the conjecture that it is also sufficient at source lines 34-36.. Background for the evenness condition that the counted lemma assumes on its input tuple x in N\textasciicircum{}\{n+1\}.
\end{itemize}\medskip
\ifdefined\hypertarget\hypertarget{sourcepointer-1}{}\fi
\noindent\textit{Source reference.} Original theorem 1.2: sources/combinatorics-fl03-matchings-in-the-hypercube-with-sp/pinned-original-source.tex, source lines 53--55. The retained context sentence at source line 64 names Theorem t:perfect when it records that the equivalence of admissibility and evenness for perfect tuples is the main theorem. The counted lemma l:first and its proof never use that theorem; the name is a navigation pointer only.\par\medskip
\begin{equation}\label{e:null}
\sum_{i=1}^{2^{n-1}} \bm{d_i} = \sum_{i=1}^{2^{n-1}}  \bm{a_i} - \bm{b_i} =   \sum_{i=1}^{2^{n-1}}  \bm{a_i} + \bm{b_i} = \sum_{\bm{v} \in \mathbb{F}_2^n} \bm{v} = \bm{0}.
\end{equation}
They conjectured that this necessary condition was also sufficient.
\begin{conjecture}[Balister, Gy{\H{o}}ri, and Schelp]\label{c:BGS}
Let $n \geq 2$ and let $\bm{d_1},\ldots, \bm{d_{2^{n-1}}}$ be elements of $\mathbb{F}_2^n$ such that $\sum_{i=1}^{2^{n-1}} \bm{d_i} = \bm{0}$. Then there exists a partition of $\mathbb{F}_2^n$ into pairs $\{\bm{a_i},\bm{b_i}\}$ such that $\bm{a_i} - \bm{b_i} = \bm{d_i}$ for each $i \in \left[2^{n-1}\right]$.
\end{conjecture}

Given a matching $M \subseteq Q^n$ we say the \emph{profile} of $M$ is the vector $\bm{x}=(x_1,\ldots, x_n) \in \mathbb{N}^n$ such that $M$ contains $x_i$ edges whose endpoints differ in the $i$th coordinate, and we are interested in which tuples $\bm{x} \in \mathbb{N}^n$ are achievable as the profile of some matching. We call such tuples \emph{admissible}. In this case, the obvious necessary condition for a tuple to be the profile of a perfect matching is that each unit vector appears an even number of times as a matching edge, i.e., each $x_i$ is even, in which case we say that $\bm{x}$ is \emph{even}, and Conjecture \ref{c:BGS} would imply that this condition is also sufficient. The aim of this short note is to verify the conjecture in this case. Let us call a tuple $\bm{x} \in \mathbb{N}^n$  \emph{perfect} if $||\bm{x}||_1 := \sum_i x_i = 2^{n-1}$.

We first note that by symmetry the ordering of the coordinates of $\bm{x}$ is irrelevant, and so we may assume from this point forwards that the coordinates appear in non-decreasing order, i.e., $x_1 \leq x_2 \leq \ldots \leq x_n$. Given $\bm{x},\bm{z} \in \mathbb{N}^{n}$ we say that \emph{$\bm{z}$ precedes $\bm{x}$}, which we write as $\bm{z} \preceq \bm{x}$, if $z_i \leq x_i$ for all $i \in [n]$. Clearly the set of admissible tuples is down-closed in the partial order $\preceq$. In particular, Theorem \ref{t:perfect} is equivalent to the statement that every even tuple is admissible. Finally, given $\bm{x} \in \mathbb{N}^{n}$ and $\bm{z} \in \mathbb{N}^{m}$ we will write $(\bm{x},\bm{z})$ for the vector in $\mathbb{N}^{n+m}$ obtained by concatenation $\bm{x}$ and $\bm{z}$.

Let us make the preceding idea formal with the following lemma. Given $\bm{x} \in \mathbb{N}^n$, let us write $o(\bm{x})$ for the number of coordinates of $\bm{x}$ except the last which are congruent to $2 \mod 4$.

\begin{lemma}\label{l:first}
Let $n\geq 2$, and let $\bm{x} \in \mathbb{N}^{n+1}$ be even with $||\bm{x}||_1 \leq 2^{n}$ and $x_{n+1} \geq 2 o(\bm{x})$. If $\bm{y} := \left(\frac{x_1}{2},\frac{x_2}{2}, \ldots, \frac{x_n}{2}\right)$, then there exists a perfect even $\bm{y'} \in \mathbb{N}^n$ such that $\bm{y} \preceq \bm{y'}$
\end{lemma}

\begin{proof}
For each $i \in [n]$ let 
\[
o_i = \begin{cases} 1 \qquad &\text{ if } x_i \equiv 2 \mod 4,\\
0 \qquad &\text{ otherwise}.\end{cases}
\]
Let $\bm{z} = \bm{y} + \bm{o}$. Note, in particular, that $\bm{y} \preceq \bm{z}$ and that by construction $\bm{z}$ is even.

Since $||\bm{o}||_1 = o(\bm{x}) \leq \frac{x_{n+1}}{2}$, it follows that 
\[
||\bm{z}|| \leq \sum_{i=1}^n \frac{x_i}{2} + o_i \leq \frac{||\bm{x}||_1}{2} \leq 2^{n-1}.
\]
Hence, since $2^{n-1}$ is even, there is some perfect even $\bm{y'} \succeq \bm{z} \succeq \bm{y}$ as claimed.
\end{proof}

\end{document}
